% Transcription — fonds Alexandre Grothendieck, shelfmark 44, batch 1 (pages 1-8).
% Established from the facsimile archives/batches/44/batch-01.pdf (a mirror of
% https://grothendieck.umontpellier.fr/44.pdf), per the skill
% transcribe-grothendieck. The folder is eight pages and this is its only batch.
% It is a single dossier, paginated 1 to 5 in his own hand across the leaves he
% wrote on: pages 1, 3 and 5 carry a "Programme de travail" of seven numbered
% theorems on intermediate Jacobians, and pages 7 and 8 carry the typescript of
% his letter to Serre of 12 August 1964, whose two leaves he numbered 4 and 5 so
% that the letter closes the same run. The rectos he wrote on were the backs of
% paper already used: page 2 is a cancelled sheet of his own on the completion of
% a graded ring, page 4 a cancelled sheet of his own on the axioms of a root
% datum, and page 6 a Bourbaki typescript in nobody's hand. Following folder 29
% batch 11, the two cancelled sheets in his hand are set at their rank with a
% note disclaiming any continuity; following folder 54, the Bourbaki typescript
% does not appear, and the gap in the numbering is the record that it was
% skipped.
% Pass: Opus 5 (claude-opus-5), 2026-08-30 — first pass, unchecked against the pages by a human.
\documentclass[11pt,a4paper]{article}
% Le préambule commun à toutes les transcriptions du fonds Grothendieck.
%
% Il tient en une idée : **l'incertitude se compose, elle ne se résout pas.**
% Une transcription qui lisse un mot douteux a détruit la seule chose pour
% laquelle on la faisait — un lecteur doit pouvoir savoir, à chaque endroit,
% ce qui était sur le feuillet et ce qui a été ajouté. Les six macros qui
% suivent sont donc l'appareil critique, et non de la décoration.
%
% Les mêmes commandes sont reconnues par `scripts/render.mjs`, qui produit la
% vue de lecture du site. Une macro ajoutée ici doit l'être aussi là-bas,
% faute de quoi le rendu échouera bruyamment — ce qui est voulu : un rendu
% silencieusement faux est pire qu'un rendu absent.

\ProvidesPackage{grothendieck}[2026/08/08 Transcriptions du fonds Grothendieck]

\usepackage{amsmath,amssymb,amsthm}
\usepackage{mathrsfs}
\usepackage{stmaryrd}
% Les diagrammes commutatifs sont partout dans ce fonds : c'est la langue
% même de la géométrie algébrique de Grothendieck, et une transcription qui
% les rendrait en prose ne transcrirait rien.
\usepackage{tikz-cd}
% Français par défaut — c'est la langue des feuillets — mais l'anglais est
% chargé aussi : la traduction et le résumé sont des documents anglais, et la
% typographie française (espaces avant les deux-points) y serait une faute.
% Ces documents basculent par \selectlanguage{english} en tête de corps.
\usepackage[english,french]{babel}
\usepackage{fontspec}
\usepackage{geometry}
\geometry{a4paper,margin=25mm,marginparwidth=28mm}
\usepackage{marginnote}
\usepackage{xcolor}
% ulem, not soul. `soul`'s \st and \ul are built on a character-by-character
% scan that XeLaTeX's font machinery breaks: the document fails with an
% undefined control sequence rather than degrading. ulem does the same two jobs
% and is Unicode-engine safe. `normalem` stops it hijacking \emph, which the
% transcriptions use for Grothendieck's own emphasis.
\usepackage[normalem]{ulem}
\usepackage{hyperref}
\hypersetup{colorlinks=true,linkcolor=black,urlcolor=black,pdfborder={0 0 0}}

% --- Métadonnées du lot -----------------------------------------------------
% Reprises telles quelles par le script de rendu, qui les affiche en tête de
% la vue de lecture. Elles fixent ce que le fichier prétend couvrir : sans
% elles, une transcription flotte sans point d'attache dans un fonds de seize
% mille feuillets.

\newcommand{\folder}[1]{\gdef\grothFolder{#1}}
\newcommand{\batch}[1]{\gdef\grothBatch{#1}}
\newcommand{\leaves}[2]{\gdef\grothFirst{#1}\gdef\grothLast{#2}}
\newcommand{\foldertitle}[1]{\gdef\grothFoldertitle{#1}}
\gdef\grothFolder{?}\gdef\grothBatch{?}\gdef\grothFirst{?}\gdef\grothLast{?}\gdef\grothFoldertitle{}

% --- Le résumé --------------------------------------------------------------
%
% Il ouvre la lecture modernisée, sous le titre « Résumé », et il est composé
% en petit. Ce n'est pas un ornement : c'est le seul passage du document qui
% ne soit pas des mathématiques, et un lecteur doit pouvoir le reconnaître
% avant de l'avoir lu — puis le sauter s'il connaît déjà le sujet, ou s'y
% arrêter s'il ne le connaît pas. Le filet qui le suit marque la reprise du
% corps.

\newenvironment{resume}
  {\small\setlength{\parskip}{0.45em}}
  {\par\medskip\hrule\medskip}

% --- Mots-clés ---------------------------------------------------------------
%
% En anglais, en clôture du résumé de la lecture modernisée : le vocabulaire
% moderne sous lequel chercher ce que les feuillets construisent. Le manifeste
% du site (`npm run manifest`) les extrait pour étiqueter la cote entière —
% cette ligne est la source unique des tags, il n'y a pas de fichier à côté.
\newcommand{\keywords}[1]{\par\smallskip\noindent
  {\footnotesize\textsc{Keywords} --- \textit{#1}}\par}

% --- L'appareil critique ----------------------------------------------------

% Le numéro de feuillet, en marge. C'est l'ancre qui permet de régler un
% désaccord sur une formule en regardant *un* feuillet plutôt qu'une chemise
% de six cents. Le site s'en sert aussi pour tourner les pages du fac-similé
% quand on fait défiler la transcription.
\newcommand{\leaf}[1]{%
  \marginnote{\footnotesize\color{gray!70}\textsf{#1}}%
  \label{leaf:#1}\ignorespaces}

% Illisible : on ne devine pas. Un mot inventé qui se lit comme les autres est
% le pire résultat possible.
\newcommand{\ill}{\textcolor{red!60!black}{[\,\ldots\,]}}

% Lecture douteuse : le mot est donné, et signalé comme incertain.
\newcommand{\uncertain}[1]{\uline{#1}}

% Ajout de l'éditeur — une lettre manquante, un mot sous-entendu.
\newcommand{\add}[1]{\textcolor{blue!50!black}{[#1]}}

% Biffé par Grothendieck lui-même. Ce qu'il a rayé fait partie du document :
% une rature dit souvent où la pensée a changé de direction.
\newcommand{\struck}[1]{\sout{#1}}

% Note du transcripteur, sur l'état matériel du feuillet.
\newcommand{\note}[1]{\par\smallskip{\small\color{gray!60!black}\textit{[#1]}}\par\smallskip}

% Note marginale de Grothendieck, distincte du corps du texte.
\newcommand{\marginal}[1]{\par\smallskip\noindent\hspace{1em}%
  {\small\color{gray!60!black}#1}\par\smallskip}

% --- Titre ------------------------------------------------------------------

\AtBeginDocument{%
  \begin{center}
    {\large\bfseries Fonds Alexandre Grothendieck --- cote n\textsuperscript{o}~\grothFolder}\\[2pt]
    {\itshape\grothFoldertitle}\\[4pt]
    {\small feuillets \grothFirst--\grothLast\ (lot \grothBatch)}\\[2pt]
    {\footnotesize\color{gray!60!black}Transcription établie d'après le fac-similé numérisé
     par l'Université de Montpellier.\\
     Les lectures incertaines sont soulignées, les illisibles notées [\,\ldots\,],
     les ajouts entre crochets.}
  \end{center}
  \vspace{1em}}

\endinput


\folder{44}
\batch{1}
\pages{1}{8}
\dating{1964}
\watermark{Édition de démonstration}
\foldertitle{Jacobiennes intermédiaires : notes manuscrites (s.d.), lettre (1964)}

\begin{document}

\section*{Programme de travail}

\note{titre de sa main, en tête du feuillet. Les feuillets 1, 3 et 5 portent sa
pagination \emph{1}, \emph{2}, \emph{3} en haut à droite, et la lettre
dactylographiée qui les suit la continue en \emph{4} et \emph{5} : le dossier
est numéroté d'un bout à l'autre par lui, et le programme et la lettre ne font
qu'une seule suite}

\page{1}

\note{sur ces trois feuillets il écrit $P^{i}(X)$ pour la jacobienne
intermédiaire ; la lettre, plus tardive, écrira $J^{i}(X)$ pour le même objet}

\textbf{Th. 1}\quad $N(\xi) \subset \mathrm{Alb}(S)$ n'est autr\add{e}
\struck{des} (ensemblistement) que le noyau de l'application naturelle

\[
  \varphi_{\xi} : \mathrm{Alb}(S) \longrightarrow
  \mathcal{Z}^{i}(X)/\mathcal{Z}^{i}_{\mathrm{alb}}(X)
\]

\note{la lettre des groupes de cycles est tantôt un $Z$ anguleux, tantôt un $\mathcal{Z}$
bouclé, sans qu'aucune distinction ne paraisse voulue ; elle est uniformisée en
$\mathcal{Z}$ ici. L'indice \emph{alb} est entouré d'un cercle aux deux endroits
où il paraît, sans doute pour marquer une notation encore provisoire}

\textbf{Th. 2}\quad La dimension de $\mathrm{Alb}(S)/N(\xi) = P_{\xi}$ est
majorée par un nombre ne dépendant que de $X$ [en fait, par
$\tfrac{1}{2}\,b^{2i-1}(X)$\,!].

Ces deux th. \uncertain{s'expliquent} \ill{}

\textbf{Th. 3}\quad $Z^{i}_{\mathrm{alg}}(X)/\mathcal{Z}^{i}_{\mathrm{alb}}(X) =
P^{i}(X)$ a une structure naturelle de V.A., et $P^{i}(X)$ et $P^{n+1-i}(X)$
sont en dualité canonique.

\marginal{« Hodge »}

\note{le mot est souligné puis barré d'un long trait ; il nomme le th. 4 comme
« Néron », « Lefschetz » et \uncertain{« Matsusaka »} nomment les th. 5, 6 et 7}

\marginal{\uncertain{Cancelé} tel quel}

\note{cette note est entourée d'un cercle et une flèche la mène au th. 4 ; la
même formule reparaît en marge de la page 7, contre l'énoncé de la lettre qui
reprend ce théorème. La lecture du premier mot n'est pas sûre}

\textbf{Th. 4}\quad Si $n = 2m-1$, alors l'autodualité canonique de $P^{m}(X)$
[donc \ill{} sur $P^{m}(X)$] est associée à une polarisation de $P^{m}(X)$.

\section*{Feuillet d'un autre document : complétion d'un anneau gradué}

\note{feuillet sans rapport avec le programme qui l'entoure, entièrement barré
de deux longues diagonales. Il est transcrit ici à son rang, sans qu'aucune
continuité avec les jacobiennes intermédiaires ne soit affirmée. Sa prose est
très rapide et ne se laisse lire qu'en partie ; les formules, elles, portent}

\page{2}

Soit $A = \coprod_{i \geq 0} A^{i}$ un anneau gradué, \ill{} degrés $\geq 0$, et
soit

\[
  \hat{A} = \prod_{i \geq 0} A_{i}
\]

le complété \ill{}. Si \ill{} noeth. [p. ex. $A$ \ill{}, $A_{1}$ de type fini
sur $A_{0}$], alors $\hat{A}$ est plat sur $A$. Pour tout $A$-module $M$, nous
posons $\hat{M} = M \otimes_{A} \hat{A}$.

Supposons \ill{} que $M$ soit \ill{} gradué de type fini, alors \ill{}

\[
  \hat{M} = \prod_{i \geq 0} M_{i}
\]

Si $M$ satisfait \ill{}, alors \ill{} une suite exacte

\[
  \cdots \longrightarrow M'' \longrightarrow M \longrightarrow M' \longrightarrow 0
\]

[$M'$ \ill{} gradué, $M''$ de type fini], et la \ill{}

\[
  \cdots \longrightarrow \hat{M}'' \longrightarrow \hat{M} \longrightarrow \hat{M}' \longrightarrow 0
\]

\[
  \hat{M}'' = \prod_{i \geq 0} M''_{i} = \prod_{i \geq 0} M_{i}
  \qquad
  \hat{M}' \overset{\sim}{\longrightarrow} M'
\]

comme on voit de suite [car $M'$ \ill{} des \ill{} de définition]. \ill{} de
sommes $\sum x_{i}$ \ill{} \struck{toujours cas} \ill{} $i \to \infty$. \ill{}
exact et fidèle.

\section*{Programme de travail (suite)}

\page{3}

\note{sa page \emph{2}}

\struck{$2i-1 + 2j-1 = 2(i+j-1) = n$}

\marginal{« Néron »}

\textbf{Th. 5}\quad Soient $X$ variété projective n.s. de dim $n$, $Z^{i},
Z'^{j} \in \mathcal{Z}^{*}(X)$ de codim $i, j$, avec \struck{codim} $i+j = n$,
\ill{} que $z = Z\cdot Y$ et $z' = Z'\cdot Y$ soient alg. équiv. à $0$ (sur $Y$)
[$Y$ section hyperplane générique d'un \uncertain{pinceau} de telles \ill{}] ;
de codim \struck{\ill{}} $i, j$, avec \struck{$\dim Y = \dim X - 1$} $i+j = n =
\dim Y + 1$, donc \struck{définissent des} \ill{} de $P^{i}(Y)$,
$P^{0}_{j}(Y)$, et par le foncteur de Néron--Tate \ill{} [l'accouplement entre
$P^{i}(Y)$ et \uncertain{$P^{j}(Y)$}] on \struck{\ill{}} définit $\langle z,
z'\rangle$. Cela posé, on a \struck{\ill{} au signe près ??}

\[
  \langle z, z'\rangle = (-1)^{m}\, Z\cdot Z'
\]

\marginal{N.B. Si $2i = 2j$ \ill{} les deux \uncertain{membres} sont nuls,
\ill{}}

\marginal{Doit \uncertain{résulter} d'un \ill{} plus \uncertain{fin} \ill{} des
\ill{}, et \ill{} un variant local \ill{}}

\note{cette seconde note est écrite dans un triangle tracé au bas du feuillet,
et n'est lisible qu'en partie}

\marginal{« Lefschetz »}

\note{barré d'un long trait, comme « Hodge » à la page 1}

\textbf{Th. 6}\quad Sous les conditions du th. 5, si \struck{$2i \leq n+1$}
$P^{i}(X) \to P^{i}(Y)$ est une \emph{isogénie} ; et si $2i = n$ le \ill{} de
$P^{i}(Y)$ est identique à l'image de $P^{i}(X)$ ; \struck{c'est égale à
$P^{i}(X)$} \struck{sauf si $2i = n = 2i'$ \ill{}} \struck{le noyau de $P^{i}(X)
\to P^{i}(Y)$ est fini} [si \struck{$2i$} $2i \geq n+1$, $P^{i}(X) \to P^{i}(Y)$
est \ill{}] \struck{$2i \leq n+1$ surjectif}

\note{l'énoncé est d'abord annoncé « Cor. », le mot étant ensuite barré et
remplacé par « Th. 6 ». Le mot \emph{isogénie} est souligné}

\section*{Feuillet d'un autre document : les axiomes d'une donnée radicielle}

\note{second feuillet sans rapport avec le programme, barré d'une longue
diagonale, et transcrit ici à son rang aux mêmes conditions que la page 2. Le
$\Delta$ y est partout repassé à l'encre lourde par-dessus le crayon}

\page{4}

Considérons les objets suivants

\begin{enumerate}
\item $(P, \Delta)$, où $P$ est un groupe abélien libre de rang fini, $\Delta$
une partie finie de $P$, satisfaisant \ill{} conditions
  \begin{itemize}
  \item[a)] $\Delta$ engendre $P$
  \item[b)] Si $a \in \Delta$, $-a \in \Delta$, mais $na \notin \Delta$ si $n
  \neq 1, -1$, $n \in \mathbb{Q}$
  \item[c)] Pour tout $a \in \Delta$, il existe une « symétrie » $S_{a}$
  \ill{} $a$ \ill{} $-a$, et telle que $S_{a}(\Delta) \subset \Delta$.
  \end{itemize}

N.B. Alors les $S_{a}$ sont \ill{} \uncertain{déterminées} \ill{}, en effet
\ill{} l'unique symétrie \ill{} suivante : \struck{\ill{}} \ill{} de la forme

\[
  (x, y) = \sum_{a \in \Delta} \langle x, a\rangle \langle y, a\rangle
  \qquad\text{sur } (P^{\mathbb{Q}})' = (P')^{\mathbb{Q}}
\]

On aura alors

\[
  S_{\alpha}x = x - \langle \gamma_{\alpha}, x\rangle\, \alpha
\]

pour un $\gamma_{\alpha} \in P'^{\mathbb{Q}}$ bien \ill{}. \ill{} que l'on a en
fait $\gamma_{\alpha} \in P'$, i.e. les

\[
  c_{\alpha\beta} = \langle \gamma_{\alpha}, \beta\rangle \ \text{ sont } \in \mathbb{Z}
\]

[provenant en fait \ill{} chaque $\alpha$ est \ill{} dans $P$, et ceci du fait
qu'il \ill{} faire partie d'un « système fondamental de racines »]

\item $(P, P', \Delta, \Delta', \varphi)$, où \struck{\ill{}} $P, P'$ sont deux
groupes abéliens libres de rang fini, \struck{\ill{}} $\Delta \subset P$,
$\Delta' \subset P'$, $\varphi$ une bijection de $\Delta$ \ill{} $\Delta'$,
\struck{\ill{}} les conditions \struck{\ill{}} étant vérifiées :
  \begin{itemize}
  \item[a)] $P^{\mathbb{Q}}$ et $P'^{\mathbb{Q}}$ sont en dualité, $\Delta$
  engendre $P$, $\Delta'$ engendre $P'$ \struck{$\alpha \in P'$ telle que
  $\varphi(\alpha) = \alpha'$}
  \item[b)] Pour tout $\alpha \in \Delta$ \struck{$\alpha' \in \Delta'$},
  \ill{}
  \[
    S_{\alpha} : x \longmapsto x - \langle \alpha', x\rangle\, \alpha
  \]
  \ill{} $\Delta$ dans $\Delta$
  \item[b')] Pour tout $\gamma \in \Delta'$, $\alpha \in \Delta$ tel que
  $\varphi(\alpha) = \gamma$,
  \[
    S'_{\alpha} : x' \longmapsto x' - \langle x', \alpha\rangle\, \alpha'
  \]
  invarie $\Delta'$
  \end{itemize}
\end{enumerate}

\marginal{Si $\alpha' = \varphi(\alpha)$, on a $\langle \alpha', \alpha\rangle =
2$}

\section*{Programme de travail (fin)}

\page{5}

\note{sa page \emph{3} ; le feuillet ne porte que cet énoncé}

\marginal{\uncertain{« Matsusaka »}}

\textbf{Th. 7}\quad L'équivalence numérique $=$ $\tau$-équivalence [cycles
algébriques sur $X$ projective non singulière].

\section*{Lettre à Serre, Bures, 12 août 1964}

\note{sa page \emph{4} : la lettre est dactylographiée et continue la
pagination du programme. Les $\ell$ et les lettres grecques manquaient à la
machine et sont encrés à la main ; ils sont restitués sans autre mention}

\page{7}

Bures le 12.8.1964

Mon cher Serre,

Je commence à avoir de faibles lumières sur la façon de faire les VA avec les
cycles algébriquement équivalents à zéro d'une variété projective non
singulière $X$ (sur un corps alg clos $k$). Si $X$ est connexe de dimension $n$,
je sais associer à chaque entier $i$ compris entre $1$ et $n$ une VA $J^{i}(X)$,
jouant le rôle d'une « jacobienne intermédiaire » au sens de Weil, correspondant
à des classes de cycles de codimension $i$ (pour une relation d'équivalence que
je ne sais pas bien identifier pour l'instant, mais qui mérite certainement
d'être appelée « équivalence d'Albanese »). Je suis moralement sûr que ce sont
« les bonnes », bien que je n'ai pas prouvé grand-chose dessus pour l'instant ;
par contre, j'ai \uncertain{dès} conjectures en pagaille. Je te signale
seulement que $J^{1} = \mathrm{Pic}^{0}$ et $J^{n} = \mathrm{Alb}^{0}$, que
$J^{i}$ et $J^{n+1-i}$ sont canoniquement duales l'une de l'autre, que $\dim
J^{i} \leq \tfrac{1}{2} b_{2i-1}$ (nombre de Betti) --- de façon plus précise
$T_{\ell}(J^{i})$ est un quotient d'un sous-module de $H^{2i-1}(X,
\underline{\mathbb{Z}}_{\ell}(i))$ (et quand on sera plus savant, on saura
prouver que c'est même un sous-module dudit), en caractéristique $0$ on peut
également remplacer la cohomologie de Weil par la cohomologie de Hodge $H^{i}(X,
\underline{\Omega}{}^{i-1}_{X})$\ldots

\note{« dès » pour « des » : la coquille est de la machine et n'est pas
corrigée}

\marginal{conjecturalement \uncertain{$\mathcal{Z}^{i}$} $\subset$
\uncertain{$H^{i}(X, \Omega^{(i)}_{X})$}}

\marginal{\struck{\ill{}}}

\note{cette seconde note marginale, de deux lignes, est barrée d'une croix et
n'est plus lisible}

Moyennant un certain nombre de choses non prouvées (qui risquent d'ailleurs
d'être vaches) le théorème de positivité de Hodge se ramène à l'énoncé suivant,
que je ne sais pas prouver, et que j'ai envie de te communiquer car il est bien
terre à terre. Moralement, il s'agit de prouver que dans le cas où $\dim X =
2m-1$, l'autodualité de $J^{m}$ (qui s'exprime par une classe de correspondance
divisorielle sur $J^{m}\times J^{m}$ \struck{\ill{}} \emph{symétrique}, donc
provenant --- du moins modulo le facteur $2$ ---

\marginal{\uncertain{Cancelé} tel quel, il faut \ill{} de la
\uncertain{primitive}\ldots}

\page{8}

\note{sa page \emph{5}}

d'un élément du groupe de Néron Severi de $J^{m}$) est \emph{positive} i.e.
l'élément en question de Néron-Severi est ample, i.e. une polarisation.

Pratiquement, cela s'explicite ainsi : Soit $T$ une variété de paramètres
connexe non singulière munie d'un point marqué $a$, $z$ une classe de cycles de
codimension $m$ sur $T\times X$ (à équivalence linéaire près, mettons), telle
que $z(a) = 0$ dans $X$, soient $p$ et $q$ les deux projections de $T\times
T\times X$ sur $T\times X$, $r$ sa projection sur $T\times T$, considérons

\[
  D = r_{*}\!\left(p^{*}(z)\cdot q^{*}(z)\right)
\]

qui est une classe de diviseurs sur $T\times T$, que nous considérons comme une
classe de correspondance divisorielle sur $T\times T$. Si $A =
\mathrm{Alb}^{0}(T)$ (NB si tu veux, tu peux supposer $T = A$ et $a$
l'origine), elle provient donc d'une classe de correspondance sur $A\times A$,
évidemment symétrique. Soit $N$ le « noyau » de cette classe (i.e. le noyau de
$A \longrightarrow$ (duale de $A$) qu'elle définit), on obtient alors une classe
de correspondance symétrique sur $J\times J$, où $J = A/N$. À prouver que cette
dernière est \emph{positive} ! Je me demande si les spécialistes « abéliens »
pourraient avoir une idée sur une telle question, peut-être Matsusaka ? Ou
toi-même ? Notes d'ailleurs que cette question te dévoile pratiquement la
méthode de construction des $J^{i}$ généraux ; si tu veux, tu peux te borner
aussi au cas où tu disposes d'une sous-variété de codimension $m-1$ $Y$ de $X$,
non singulière si tu y tiens, \struck{\ill{}} et où $T = \mathrm{Pic}^{0}(Y)$,
considéré comme paramétrant les diviseurs alg équiv à zéro de $Y$, mais
considérés comme classes de cycles \struck{\ill{}} de codimension $m$ de $X$.

\note{« connexe » et « aussi » sont ajoutés en interligne à la machine}

Merci pour la copie de la lettre à Ogg !

Bien à toi

\end{document}
